% SPEC-DB-001 --- Schema Contract
% Shared preamble for all spec documents.
% Include from each spec: % Shared preamble for all spec documents.
% Include from each spec: % Shared preamble for all spec documents.
% Include from each spec: \input{../preamble.tex} (adjust depth).

\documentclass[11pt,a4paper]{article}

\usepackage[utf8]{inputenc}
\usepackage[T1]{fontenc}
\usepackage{lmodern}
\usepackage[a4paper,margin=2.3cm]{geometry}
\usepackage{parskip}
\usepackage{microtype}
\usepackage[dvipsnames]{xcolor}
\usepackage{enumitem}
\usepackage{listings}
\usepackage{tcolorbox}
\usepackage{titlesec}
\usepackage{fancyhdr}
\usepackage{array}
\usepackage{longtable}
\usepackage{booktabs}
\usepackage{amsmath}
\usepackage{amssymb}
\usepackage{hyperref}
\usepackage{cleveref}

\tcbuselibrary{skins,breakable}

\definecolor{specBlue}{RGB}{30,64,132}
\definecolor{specGray}{RGB}{90,90,90}
\definecolor{specAccent}{RGB}{198,40,40}
\definecolor{specGreen}{RGB}{30,110,60}

\hypersetup{
  colorlinks=true,
  linkcolor=specBlue,
  urlcolor=specBlue,
  citecolor=specBlue,
  pdfborder={0 0 0}
}

\titleformat{\section}{\Large\bfseries\color{specBlue}}{\thesection}{0.75em}{}
\titleformat{\subsection}{\large\bfseries\color{specBlue!80}}{\thesubsection}{0.6em}{}

\makeatletter
\newcommand{\specID}[1]{\def\@specID{#1}}
\newcommand{\specTitle}[1]{\def\@specTitle{#1}}
\newcommand{\specStatus}[1]{\def\@specStatus{#1}}
\newcommand{\specVersion}[1]{\def\@specVersion{#1}}
\newcommand{\specOwner}[1]{\def\@specOwner{#1}}
\newcommand{\specTraces}[1]{\def\@specTraces{#1}}

\specID{SPEC-XXX}
\specTitle{Untitled Spec}
\specStatus{DRAFT}
\specVersion{0.1}
\specOwner{Jeremie Nunez}
\specTraces{}

\newcommand{\makespecheader}{%
  \begin{center}
    {\LARGE\bfseries\color{specBlue}\@specTitle}\\[0.4em]
    {\small\color{specGray}\texttt{\@specID} \quad v\@specVersion \quad \textsc{\@specStatus} \quad \@specOwner}
  \end{center}
  \vspace{0.4em}
  \hrule
  \vspace{0.8em}
}

\pagestyle{fancy}
\fancyhf{}
\fancyhead[L]{\small\color{specGray}\@specID}
\fancyhead[R]{\small\color{specGray}\@specTitle}
\fancyfoot[C]{\small\color{specGray}\thepage}
\renewcommand{\headrulewidth}{0pt}
\makeatother

\newtcolorbox{invariant}{
  colback=specBlue!5, colframe=specBlue!75,
  title=Invariant, fonttitle=\bfseries, breakable,
  boxrule=0.5pt, arc=2pt
}

\newtcolorbox{scenario}[1]{
  colback=white, colframe=specBlue!50,
  title={Scenario: #1}, fonttitle=\bfseries, breakable,
  boxrule=0.5pt, arc=2pt
}

\newtcolorbox{ac}[1]{
  colback=specAccent!5, colframe=specAccent!60,
  title={Acceptance Criterion #1}, fonttitle=\bfseries, breakable,
  boxrule=0.5pt, arc=2pt
}

\newtcolorbox{nongoal}{
  colback=specGray!8, colframe=specGray!60,
  title=Non-Goal, fonttitle=\bfseries, breakable,
  boxrule=0.5pt, arc=2pt
}

\lstdefinestyle{codestyle}{
  basicstyle=\ttfamily\small,
  breaklines=true,
  showstringspaces=false,
  frame=single,
  framesep=4pt,
  rulecolor=\color{specBlue!30},
  backgroundcolor=\color{specBlue!2},
  xleftmargin=0.5em,
  xrightmargin=0.5em,
  keywordstyle=\color{specBlue}\bfseries,
  commentstyle=\color{specGray}\itshape,
  stringstyle=\color{specGreen}
}
\lstset{
  inputencoding=utf8,
  extendedchars=true,
  literate=
    {—}{{--}}1
    {–}{{-}}1
    {…}{{...}}1
    {’}{{'}}1
    {‘}{{'}}1
    {“}{{``}}1
    {”}{{''}}1
    {é}{{\'e}}1
    {è}{{\`e}}1
    {ê}{{\^e}}1
    {à}{{\`a}}1
    {â}{{\^a}}1
    {ç}{{\c{c}}}1
    {ô}{{\^o}}1
    {→}{{$\to$}}1
    {←}{{$\leftarrow$}}1
    {×}{{$\times$}}1
    {≥}{{$\geq$}}1
    {≤}{{$\leq$}}1
    {≠}{{$\neq$}}1
    {∈}{{$\in$}}1
    {⌈}{{$\lceil$}}1
    {⌉}{{$\rceil$}}1
    {∞}{{$\infty$}}1
    {α}{{$\alpha$}}1
    {β}{{$\beta$}}1
}
\lstdefinelanguage{TypeScript}{
  keywords={const,let,var,function,return,if,else,for,while,class,interface,type,extends,implements,import,export,from,as,async,await,new,this,true,false,null,undefined,void,any,unknown,never,string,number,boolean,readonly,private,public,protected,enum},
  sensitive=true,
  comment=[l]{//},
  morecomment=[s]{/*}{*/},
  morestring=[b]",
  morestring=[b]',
  morestring=[b]`
}
\lstset{style=codestyle}

\newcommand{\Given}{\textbf{\color{specBlue}Given} }
\newcommand{\When}{\textbf{\color{specBlue}When} }
\newcommand{\Then}{\textbf{\color{specBlue}Then} }
\providecommand{\And}{}
\renewcommand{\And}{\textbf{\color{specBlue}And} }
 (adjust depth).

\documentclass[11pt,a4paper]{article}

\usepackage[utf8]{inputenc}
\usepackage[T1]{fontenc}
\usepackage{lmodern}
\usepackage[a4paper,margin=2.3cm]{geometry}
\usepackage{parskip}
\usepackage{microtype}
\usepackage[dvipsnames]{xcolor}
\usepackage{enumitem}
\usepackage{listings}
\usepackage{tcolorbox}
\usepackage{titlesec}
\usepackage{fancyhdr}
\usepackage{array}
\usepackage{longtable}
\usepackage{booktabs}
\usepackage{amsmath}
\usepackage{amssymb}
\usepackage{hyperref}
\usepackage{cleveref}

\tcbuselibrary{skins,breakable}

\definecolor{specBlue}{RGB}{30,64,132}
\definecolor{specGray}{RGB}{90,90,90}
\definecolor{specAccent}{RGB}{198,40,40}
\definecolor{specGreen}{RGB}{30,110,60}

\hypersetup{
  colorlinks=true,
  linkcolor=specBlue,
  urlcolor=specBlue,
  citecolor=specBlue,
  pdfborder={0 0 0}
}

\titleformat{\section}{\Large\bfseries\color{specBlue}}{\thesection}{0.75em}{}
\titleformat{\subsection}{\large\bfseries\color{specBlue!80}}{\thesubsection}{0.6em}{}

\makeatletter
\newcommand{\specID}[1]{\def\@specID{#1}}
\newcommand{\specTitle}[1]{\def\@specTitle{#1}}
\newcommand{\specStatus}[1]{\def\@specStatus{#1}}
\newcommand{\specVersion}[1]{\def\@specVersion{#1}}
\newcommand{\specOwner}[1]{\def\@specOwner{#1}}
\newcommand{\specTraces}[1]{\def\@specTraces{#1}}

\specID{SPEC-XXX}
\specTitle{Untitled Spec}
\specStatus{DRAFT}
\specVersion{0.1}
\specOwner{Jeremie Nunez}
\specTraces{}

\newcommand{\makespecheader}{%
  \begin{center}
    {\LARGE\bfseries\color{specBlue}\@specTitle}\\[0.4em]
    {\small\color{specGray}\texttt{\@specID} \quad v\@specVersion \quad \textsc{\@specStatus} \quad \@specOwner}
  \end{center}
  \vspace{0.4em}
  \hrule
  \vspace{0.8em}
}

\pagestyle{fancy}
\fancyhf{}
\fancyhead[L]{\small\color{specGray}\@specID}
\fancyhead[R]{\small\color{specGray}\@specTitle}
\fancyfoot[C]{\small\color{specGray}\thepage}
\renewcommand{\headrulewidth}{0pt}
\makeatother

\newtcolorbox{invariant}{
  colback=specBlue!5, colframe=specBlue!75,
  title=Invariant, fonttitle=\bfseries, breakable,
  boxrule=0.5pt, arc=2pt
}

\newtcolorbox{scenario}[1]{
  colback=white, colframe=specBlue!50,
  title={Scenario: #1}, fonttitle=\bfseries, breakable,
  boxrule=0.5pt, arc=2pt
}

\newtcolorbox{ac}[1]{
  colback=specAccent!5, colframe=specAccent!60,
  title={Acceptance Criterion #1}, fonttitle=\bfseries, breakable,
  boxrule=0.5pt, arc=2pt
}

\newtcolorbox{nongoal}{
  colback=specGray!8, colframe=specGray!60,
  title=Non-Goal, fonttitle=\bfseries, breakable,
  boxrule=0.5pt, arc=2pt
}

\lstdefinestyle{codestyle}{
  basicstyle=\ttfamily\small,
  breaklines=true,
  showstringspaces=false,
  frame=single,
  framesep=4pt,
  rulecolor=\color{specBlue!30},
  backgroundcolor=\color{specBlue!2},
  xleftmargin=0.5em,
  xrightmargin=0.5em,
  keywordstyle=\color{specBlue}\bfseries,
  commentstyle=\color{specGray}\itshape,
  stringstyle=\color{specGreen}
}
\lstset{
  inputencoding=utf8,
  extendedchars=true,
  literate=
    {—}{{--}}1
    {–}{{-}}1
    {…}{{...}}1
    {’}{{'}}1
    {‘}{{'}}1
    {“}{{``}}1
    {”}{{''}}1
    {é}{{\'e}}1
    {è}{{\`e}}1
    {ê}{{\^e}}1
    {à}{{\`a}}1
    {â}{{\^a}}1
    {ç}{{\c{c}}}1
    {ô}{{\^o}}1
    {→}{{$\to$}}1
    {←}{{$\leftarrow$}}1
    {×}{{$\times$}}1
    {≥}{{$\geq$}}1
    {≤}{{$\leq$}}1
    {≠}{{$\neq$}}1
    {∈}{{$\in$}}1
    {⌈}{{$\lceil$}}1
    {⌉}{{$\rceil$}}1
    {∞}{{$\infty$}}1
    {α}{{$\alpha$}}1
    {β}{{$\beta$}}1
}
\lstdefinelanguage{TypeScript}{
  keywords={const,let,var,function,return,if,else,for,while,class,interface,type,extends,implements,import,export,from,as,async,await,new,this,true,false,null,undefined,void,any,unknown,never,string,number,boolean,readonly,private,public,protected,enum},
  sensitive=true,
  comment=[l]{//},
  morecomment=[s]{/*}{*/},
  morestring=[b]",
  morestring=[b]',
  morestring=[b]`
}
\lstset{style=codestyle}

\newcommand{\Given}{\textbf{\color{specBlue}Given} }
\newcommand{\When}{\textbf{\color{specBlue}When} }
\newcommand{\Then}{\textbf{\color{specBlue}Then} }
\providecommand{\And}{}
\renewcommand{\And}{\textbf{\color{specBlue}And} }
 (adjust depth).

\documentclass[11pt,a4paper]{article}

\usepackage[utf8]{inputenc}
\usepackage[T1]{fontenc}
\usepackage{lmodern}
\usepackage[a4paper,margin=2.3cm]{geometry}
\usepackage{parskip}
\usepackage{microtype}
\usepackage[dvipsnames]{xcolor}
\usepackage{enumitem}
\usepackage{listings}
\usepackage{tcolorbox}
\usepackage{titlesec}
\usepackage{fancyhdr}
\usepackage{array}
\usepackage{longtable}
\usepackage{booktabs}
\usepackage{amsmath}
\usepackage{amssymb}
\usepackage{hyperref}
\usepackage{cleveref}

\tcbuselibrary{skins,breakable}

\definecolor{specBlue}{RGB}{30,64,132}
\definecolor{specGray}{RGB}{90,90,90}
\definecolor{specAccent}{RGB}{198,40,40}
\definecolor{specGreen}{RGB}{30,110,60}

\hypersetup{
  colorlinks=true,
  linkcolor=specBlue,
  urlcolor=specBlue,
  citecolor=specBlue,
  pdfborder={0 0 0}
}

\titleformat{\section}{\Large\bfseries\color{specBlue}}{\thesection}{0.75em}{}
\titleformat{\subsection}{\large\bfseries\color{specBlue!80}}{\thesubsection}{0.6em}{}

\makeatletter
\newcommand{\specID}[1]{\def\@specID{#1}}
\newcommand{\specTitle}[1]{\def\@specTitle{#1}}
\newcommand{\specStatus}[1]{\def\@specStatus{#1}}
\newcommand{\specVersion}[1]{\def\@specVersion{#1}}
\newcommand{\specOwner}[1]{\def\@specOwner{#1}}
\newcommand{\specTraces}[1]{\def\@specTraces{#1}}

\specID{SPEC-XXX}
\specTitle{Untitled Spec}
\specStatus{DRAFT}
\specVersion{0.1}
\specOwner{Jeremie Nunez}
\specTraces{}

\newcommand{\makespecheader}{%
  \begin{center}
    {\LARGE\bfseries\color{specBlue}\@specTitle}\\[0.4em]
    {\small\color{specGray}\texttt{\@specID} \quad v\@specVersion \quad \textsc{\@specStatus} \quad \@specOwner}
  \end{center}
  \vspace{0.4em}
  \hrule
  \vspace{0.8em}
}

\pagestyle{fancy}
\fancyhf{}
\fancyhead[L]{\small\color{specGray}\@specID}
\fancyhead[R]{\small\color{specGray}\@specTitle}
\fancyfoot[C]{\small\color{specGray}\thepage}
\renewcommand{\headrulewidth}{0pt}
\makeatother

\newtcolorbox{invariant}{
  colback=specBlue!5, colframe=specBlue!75,
  title=Invariant, fonttitle=\bfseries, breakable,
  boxrule=0.5pt, arc=2pt
}

\newtcolorbox{scenario}[1]{
  colback=white, colframe=specBlue!50,
  title={Scenario: #1}, fonttitle=\bfseries, breakable,
  boxrule=0.5pt, arc=2pt
}

\newtcolorbox{ac}[1]{
  colback=specAccent!5, colframe=specAccent!60,
  title={Acceptance Criterion #1}, fonttitle=\bfseries, breakable,
  boxrule=0.5pt, arc=2pt
}

\newtcolorbox{nongoal}{
  colback=specGray!8, colframe=specGray!60,
  title=Non-Goal, fonttitle=\bfseries, breakable,
  boxrule=0.5pt, arc=2pt
}

\lstdefinestyle{codestyle}{
  basicstyle=\ttfamily\small,
  breaklines=true,
  showstringspaces=false,
  frame=single,
  framesep=4pt,
  rulecolor=\color{specBlue!30},
  backgroundcolor=\color{specBlue!2},
  xleftmargin=0.5em,
  xrightmargin=0.5em,
  keywordstyle=\color{specBlue}\bfseries,
  commentstyle=\color{specGray}\itshape,
  stringstyle=\color{specGreen}
}
\lstset{
  inputencoding=utf8,
  extendedchars=true,
  literate=
    {—}{{--}}1
    {–}{{-}}1
    {…}{{...}}1
    {’}{{'}}1
    {‘}{{'}}1
    {“}{{``}}1
    {”}{{''}}1
    {é}{{\'e}}1
    {è}{{\`e}}1
    {ê}{{\^e}}1
    {à}{{\`a}}1
    {â}{{\^a}}1
    {ç}{{\c{c}}}1
    {ô}{{\^o}}1
    {→}{{$\to$}}1
    {←}{{$\leftarrow$}}1
    {×}{{$\times$}}1
    {≥}{{$\geq$}}1
    {≤}{{$\leq$}}1
    {≠}{{$\neq$}}1
    {∈}{{$\in$}}1
    {⌈}{{$\lceil$}}1
    {⌉}{{$\rceil$}}1
    {∞}{{$\infty$}}1
    {α}{{$\alpha$}}1
    {β}{{$\beta$}}1
}
\lstdefinelanguage{TypeScript}{
  keywords={const,let,var,function,return,if,else,for,while,class,interface,type,extends,implements,import,export,from,as,async,await,new,this,true,false,null,undefined,void,any,unknown,never,string,number,boolean,readonly,private,public,protected,enum},
  sensitive=true,
  comment=[l]{//},
  morecomment=[s]{/*}{*/},
  morestring=[b]",
  morestring=[b]',
  morestring=[b]`
}
\lstset{style=codestyle}

\newcommand{\Given}{\textbf{\color{specBlue}Given} }
\newcommand{\When}{\textbf{\color{specBlue}When} }
\newcommand{\Then}{\textbf{\color{specBlue}Then} }
\providecommand{\And}{}
\renewcommand{\And}{\textbf{\color{specBlue}And} }


\specID{SPEC-DB-001}
\specTitle{Schema Contract --- typed entity shapes and relations}
\specStatus{DRAFT}
\specVersion{0.1}
\specOwner{Jeremie Nunez}

\begin{document}
\makespecheader

\section{Context}
The \texttt{@interview/db-schema} package is the single source of truth for the
relational model consumed by \texttt{@interview/thalamus} (research orchestration)
and \texttt{@interview/sweep} (audit pipeline). It is a thin Drizzle ORM layer on
PostgreSQL: table definitions, column types, indexes, and \texttt{\$inferSelect /
\$inferInsert} types. Every downstream module imports entity shapes from here;
drift between database and types is a class of bug this spec rules out.

\section{Goals}
\begin{itemize}[leftmargin=1.2em]
  \item Publish a stable, typed contract for every table: user, article,
        research-cycle, research-finding, research-edge, exploration-log, and
        the domain entity cluster (entity, category, attribute, relation,
        reference-entity).
  \item Encode relations as foreign-key references at the schema level so the
        ORM can enforce join keys without raw SQL.
  \item Declare indexes that support the access patterns used by repositories
        (lookup by entity name/type, lookup by cycle).
  \item Expose PostgreSQL-specific features (\texttt{jsonb}, embedding columns,
        timestamp with time zone) through typed columns only.
\end{itemize}

\section{Non-Goals}
\begin{nongoal}
This spec does not cover migrations (order, reversibility, data backfill),
row-level security, partitioning, or connection pooling. It does not prescribe
a vector index implementation (HNSW vs IVFFlat) --- embeddings are stored as
\texttt{jsonb} payloads at this stage and similarity search lives in a separate
spec. Business validation rules (Zod schemas, DTO mapping) are out of scope.
\end{nongoal}

\section{Invariants}
\begin{invariant}
Every table exposes \texttt{\$inferSelect} and, where inserts are performed,
\texttt{\$inferInsert}. Consumers never hand-roll row types; they import them
from \texttt{@interview/db-schema}. Column nullability in TypeScript matches
\texttt{notNull()} in the schema exactly.
\end{invariant}

\begin{invariant}
Primary keys are \texttt{bigserial} exposed as \texttt{bigint}. Foreign keys
use \texttt{.references(() => table.id)} so Drizzle can participate in typed
joins. Timestamps are \texttt{timestamp with time zone} with \texttt{defaultNow()}
for \texttt{createdAt}.
\end{invariant}

\begin{invariant}
Free-form payloads (plans, results, metadata, embeddings, sensory descriptors)
use \texttt{jsonb}. Repositories are responsible for narrowing these at read
time; the schema does not leak \texttt{any}.
\end{invariant}

\section{Interface}

The public surface is re-exported from \texttt{packages/db-schema/src/index.ts}:
the \texttt{Database} alias and every table + inferred type.

\begin{lstlisting}[language=TypeScript]
// packages/db-schema/src/index.ts
import type { NodePgDatabase } from "drizzle-orm/node-postgres";
import * as schema from "./schema";

export type Database = NodePgDatabase<typeof schema>;
export * from "./schema";

// Representative entity shapes (re-exported)
export type User            = typeof user.$inferSelect;
export type Article         = typeof article.$inferSelect;
export type NewArticle      = typeof article.$inferInsert;
export type ResearchCycle   = typeof researchCycle.$inferSelect;
export type ResearchFinding = typeof researchFinding.$inferSelect;
export type ResearchEdge    = typeof researchEdge.$inferSelect;
export type ExplorationLog  = typeof explorationLog.$inferSelect;
\end{lstlisting}

\subsection{Entity groups}

\begin{center}
\small
\begin{tabular}{@{}lll@{}}
\toprule
Group & Tables & Purpose \\
\midrule
Identity   & \texttt{user}                                    & Actor, role, tier. \\
Content    & \texttt{article}                                 & Authored records, FK to \texttt{user}. \\
Research   & \texttt{research\_cycle}, \texttt{research\_finding}, \texttt{research\_edge} & Orchestration state + knowledge graph. \\
Audit      & \texttt{exploration\_log}                        & Topic-scoped run traces. \\
Reference  & entity cluster (\texttt{entity}, \texttt{category}, \texttt{attribute}, \texttt{relation}, \texttt{reference\_entity}) & Domain reference data consumed by research and audit. \\
\bottomrule
\end{tabular}
\end{center}

\subsection{Indexes and access shapes}
\begin{itemize}[leftmargin=1.2em]
  \item \texttt{idx\_research\_finding\_entity} on \texttt{(entity\_name, entity\_type)}:
        supports "find everything known about entity X".
  \item \texttt{idx\_research\_finding\_cycle} on \texttt{(cycle\_id)}:
        supports "stream findings for an active cycle".
  \item Unique indexes: \texttt{user.email}, \texttt{article.slug}.
\end{itemize}

\subsection{Columns of note}
\begin{itemize}[leftmargin=1.2em]
  \item \texttt{research\_finding.embedding}: \texttt{jsonb} payload reserved for
        a pgvector column in a later revision. Contract today: an array of
        \texttt{number}, length fixed per model.
  \item \texttt{research\_finding.confidence}, \texttt{research\_edge.confidence}:
        \texttt{real}, conventionally in \([0, 1]\). Not enforced by a check
        constraint --- the service layer clamps.
  \item Status-like columns (\texttt{research\_cycle.status},
        \texttt{article.status}, \texttt{exploration\_log.status}) are \texttt{text}
        today. Promotion to a \texttt{pgEnum} is tracked in Open Questions.
\end{itemize}

\section{Scenarios}

\begin{scenario}{Consumer imports an entity type}
\Given a downstream package depends on \texttt{@interview/db-schema} \\
\When it writes \texttt{import type \{ ResearchFinding \} from "@interview/db-schema"} \\
\Then the type reflects every column of \texttt{research\_finding} with correct
nullability \\
\And renaming a column in the schema is a type error at the import site.
\end{scenario}

\begin{scenario}{Insert a new finding}
\Given a \texttt{Database} handle \\
\When a service calls \texttt{db.insert(researchFinding).values(x)} with
\texttt{x : NewResearchFinding} \\
\Then TypeScript enforces required columns (\texttt{cortex},
\texttt{entityType}, \texttt{entityName}, \texttt{findingType}) \\
\And optional columns (\texttt{content}, \texttt{confidence}, \texttt{embedding},
\texttt{metadata}) may be omitted.
\end{scenario}

\begin{scenario}{Graph traversal by entity}
\Given at least one finding exists for \texttt{(entityName, entityType) = ("X","item")} \\
\When a repository queries \texttt{research\_finding} filtered by that tuple \\
\Then the planner uses \texttt{idx\_research\_finding\_entity} \\
\And latency stays below the threshold asserted by the repo-level test.
\end{scenario}

\begin{scenario}{Foreign key integrity}
\Given an \texttt{article} row with \texttt{authorId = 42} \\
\When there is no \texttt{user} with \texttt{id = 42} \\
\Then the database rejects the insert \\
\And the typed \texttt{Article} reflects \texttt{authorId: bigint | null}.
\end{scenario}

\begin{scenario}{JSONB payload is opaque at schema level}
\Given \texttt{research\_cycle.plan} is \texttt{jsonb} \\
\When a consumer reads a \texttt{ResearchCycle} \\
\Then \texttt{plan} is typed as \texttt{unknown} (or the schema-declared narrow
type) --- never \texttt{any} \\
\And narrowing is performed by the service, not the schema.
\end{scenario}

\section{Acceptance Criteria}

\begin{ac}{AC-1}
\texttt{pnpm --filter @interview/db-schema typecheck} passes with
\texttt{strict: true} and \texttt{noUncheckedIndexedAccess: true}.
\end{ac}

\begin{ac}{AC-2}
Every table declared under \texttt{src/schema/} is re-exported from
\texttt{src/schema/index.ts} and transitively from \texttt{src/index.ts}. A test
asserts parity between the file system and the export list.
\end{ac}

\begin{ac}{AC-3}
Every \texttt{notNull()} column maps to a non-optional, non-nullable field in
\texttt{\$inferSelect}. A type-level test (\texttt{expectType}) locks this in
for at least \texttt{user}, \texttt{research\_cycle}, \texttt{research\_finding},
\texttt{article}.
\end{ac}

\begin{ac}{AC-4}
Foreign keys declared in the schema match the referenced table's primary key
type (\texttt{bigint}). A compile-time fixture verifies
\texttt{article.authorId} references \texttt{user.id} and
\texttt{research\_finding.cycleId} references \texttt{researchCycle.id}.
\end{ac}

\begin{ac}{AC-5}
Indexes \texttt{idx\_research\_finding\_entity} and
\texttt{idx\_research\_finding\_cycle} exist in the generated SQL and are
covered by an integration test that runs \texttt{EXPLAIN} and asserts an
index scan.
\end{ac}

\begin{ac}{AC-6}
Searching the schema sources for \texttt{: any} returns zero matches.
Searching for \texttt{jsonb(} returns only the documented payload columns.
\end{ac}

\section{Traceability}

\begin{center}
\small
\begin{tabular}{@{}lll@{}}
\toprule
AC & Test file (planned) & Test name \\
\midrule
AC-1 & \texttt{packages/db-schema/tests/typecheck.spec.ts}       & \texttt{schema typechecks under strict} \\
AC-2 & \texttt{packages/db-schema/tests/exports.spec.ts}         & \texttt{every table file is re-exported} \\
AC-3 & \texttt{packages/db-schema/tests/nullability.spec.ts}     & \texttt{notNull columns are non-nullable in \$inferSelect} \\
AC-4 & \texttt{packages/db-schema/tests/fk-types.spec.ts}        & \texttt{foreign keys match referenced pk types} \\
AC-5 & \texttt{packages/db-schema/tests/indexes.int.spec.ts}     & \texttt{finding lookups use expected indexes} \\
AC-6 & \texttt{packages/db-schema/tests/no-any.spec.ts}          & \texttt{schema sources contain no \textbackslash any} \\
\bottomrule
\end{tabular}
\end{center}

\section{Open Questions}
\begin{itemize}[leftmargin=1.2em]
  \item Promote status-like \texttt{text} columns to \texttt{pgEnum} now or in a
        follow-up migration?
  \item Swap \texttt{research\_finding.embedding} from \texttt{jsonb} to a
        native \texttt{vector(n)} column --- which \(n\), which index?
  \item Should \texttt{research\_edge} carry a composite unique key on
        \texttt{(fromName, fromType, toName, toType, relation)} to deduplicate
        graph edges at the database level?
  \item Do we need soft-delete (\texttt{deletedAt}) on \texttt{article} and
        \texttt{user} for audit alignment with \texttt{@interview/sweep}?
\end{itemize}

\end{document}
